\begin{solapartat}
\punts{0,25} Ens diuen que la flauta travessera és un tub obert pels dos extrems. Les longituds ressonants seran: $\lambda_n = \frac{2\,L}{n}$ i utilitzant $\lambda\cdot f = v$; les freqüències ressonants seran $f_n = n\,\frac{v}{2\,L} = n\,\frac{343}{2\cdot 0,67} = n\cdot 256\ \mathrm{Hz}$. Per tant, el primer harmònic serà $f_1 = 256\ \mathrm{Hz}$.

\punts{0,25} Igualment, per al tercer harmònic: $f_3 = 3\cdot 256\ Hz = 768\ \mathrm{Hz}$.

\punts{0,25} Tub amb els dos extrems oberts; té els ventres als extrems i la distància entre ventres de l'ona estacionària és mitja longitud d'ona, posició $x_V = n\,\frac{\lambda}{2}$. Els nodes entre ventres a mitja distància, $x_N = (2n+1)\,\frac{\lambda}{4}$.

\punts{0,25} Per al primer harmònic: $\frac{\lambda_1}{2} = 0,67\ \mathrm{m}$; per tant: primer ventre: $x_{V1} = 0$; segon ventre: $x_{V2} = 0,67\ \mathrm{m}$, Node: $x_{N1} = 0,335\ \mathrm{m}$.

\punts{0,25} Per al tercer harmònic: $\frac{\lambda_3}{2} = 0,223\ \mathrm{m}$; per tant:

\begin{tabular}{@{}l}
Primer ventre: $x_{V1} = 0$\\
Segon ventre: $x_{V2} = 0,223\ \mathrm{m}$\\
Tercer ventre: $x_{V3} = 0,447\ \mathrm{m}$\\
Quart ventre: $x_{V4} = 0,67\ \mathrm{m}$\\
Primer node: $x_{N1} = 0,112\ \mathrm{m}$\\
Segon node: $x_{N2} = 0,335\ \mathrm{m}$\\
Tercer node: $x_{N3} = 0,558\ \mathrm{m}$
\end{tabular}

\figura[0.75]{sol_fig1.png}
\end{solapartat}

\begin{solapartat}
\punts{0,4} En el cas extrem tancat - extrem obert, la longitud d'ona de les ones estacionàries compleix $\lambda_n = \frac{4\,L}{n}$ i les seves freqüències corresponents són $f_n = \frac{v}{\lambda_n} = n\,\frac{v}{4\,L}$ amb $n = 1, 3, 5$.

Per tant, el primer mode de vibració correspon a $n = 1$ i té longitud d'ona: $\lambda_1 = \frac{4\cdot 0,67}{1} = 2,68\ \mathrm{m}$ i freqüència $f_{n1} = \frac{343}{2,68} = 128\ \mathrm{Hz}$.

\punts{0,35} El segon mode de vibració correspon a $n = 3$ i té longitud d'ona: $\lambda_3 = \frac{4\cdot 0,67}{3} = 0,893\ \mathrm{m}$ i freqüència $f_3 = \frac{343}{0,893} = 384\ \mathrm{Hz}$.

\punts{0,25} Dibuix del primer mode de vibració, harmònic fonamental.

\punts{0,25} Dibuix del segon mode de vibració.

\figura[0.6]{sol_fig2.png}
\end{solapartat}
